% =====================================================================
%  2bat-ud4-10.tex  —  SESSIÓ 10: Optimització: casos socials
%  (sense preàmbul: s'inclou des de main.tex)
% =====================================================================
\horatitol{Sessió 10 — Optimització: casos socials}%
{Continuació de l'Activitat 3: més problemes d'optimització, amb plantejaments una mica més exigents i contextos variats.}

\subsection{Idea clau}

A la sessió 9 vau aprendre el \textbf{mètode dels 5 passos} per optimitzar. Avui el
\textbf{posem en pràctica} amb casos més rics: problemes on construir la funció demana
una mica més de feina (una condició que cal aïllar, un volum, una superfície dividida).
El procediment és el mateix; el que millorem és la \textbf{traducció} de l'enunciat.

\begin{idea}
\textbf{Recordatori del mètode (5 passos)}
\begin{enumerate}[leftmargin=1.6em, itemsep=1pt]
  \item Identifica l'\textbf{objectiu} (què es fa màxim o mínim) i la \textbf{variable}.
  \item Escriu la \textbf{funció}; si hi ha una condició, fes-la servir per quedar-te amb
        una sola variable.
  \item \textbf{Deriva} i resol $f'(x)=0$.
  \item \textbf{Classifica} el punt crític (signe de $f'$).
  \item \textbf{Interpreta} i comprova que la solució tingui sentit en el context.
\end{enumerate}
\textbf{Consell:} sempre que puguis, escriu primer la \emph{condició} (el total, el
perímetre, el material disponible) i aïlla-hi una variable abans de muntar la funció.
\end{idea}

\subsection{Exemples}

\begin{exemple}[la satisfacció d'un casal segons la participació]
La satisfacció mitjana en una activitat de casal (en una escala interna) segons el
nombre de participants $x$ (en desenes) es modelitza amb $S(x)=-\tfrac{1}{2}x^2+30x$.
Quina participació la fa màxima?
\[
S'(x)=-x+30=0 \Rightarrow x=30.
\]
$S'$ passa de $+$ ($S'(29)=1$) a $-$ ($S'(31)=-1$): és un \textbf{màxim}. La satisfacció
és màxima amb $x=30$ desenes, és a dir $300$ participants, amb
$S(30)=-450+900=450$. Més enllà d'aquesta xifra, la massificació en fa baixar la
qualitat.
\end{exemple}

\begin{exemple}[una caixa oberta de cartolina]
A partir d'una cartolina quadrada de $12$ cm de costat es fa una caixa \textbf{sense
tapa} retallant un quadrat de costat $x$ a cada cantonada i aixecant les vores. Quin
retall dóna el \textbf{volum màxim}?

\medskip
\noindent\textbf{Funció:} la base de la caixa és un quadrat de costat $12-2x$ i l'altura
és $x$, així que
\[
V(x)=x\,(12-2x)^2=4x^3-48x^2+144x.
\]
\noindent\textbf{Derivem:} $V'(x)=12x^2-96x+144=12(x-2)(x-6)$, que s'anul·la a $x=2$ i
$x=6$. Com que el costat de la base ha de ser positiu ($12-2x>0\Rightarrow x<6$), només
té sentit $x=2$. Allà $V'$ passa de $+$ a $-$: \textbf{màxim}. El retall òptim és $x=2$
cm i el volum màxim $V(2)=2\cdot 8^2=128\ \text{cm}^3$.
\end{exemple}

\subsection{Exercicis resolts}

\begin{exresolt}[una tanca amb divisió interior]
Una entitat vol tancar un recinte rectangular i, a més, dividir-lo en dues parcel·les
iguals amb una tanca interior \textbf{paral·lela} a un dels costats. Disposa de $120$ m
de tanca en total. Quines mides donen la màxima superfície?
\solucio
\textbf{Variable i condició:} diem $x$ al costat perpendicular a la divisió i $y$ a
l'altre. La tanca interior fa que hi hagi \textbf{tres} segments de longitud $x$ i
\textbf{dos} de longitud $y$:
\[
3x+2y=120 \ \Rightarrow\ y=\frac{120-2x}{3}.
\]
\textbf{Funció (àrea):}
\[
A(x)=x\cdot y=x\cdot\frac{120-2x}{3}=\frac{120x-2x^2}{3}.
\]
\textbf{Derivem:} $A'(x)=\dfrac{120-4x}{3}=0 \Rightarrow x=30$. $A'$ passa de $+$ a $-$:
\textbf{màxim}. Aleshores $y=\dfrac{120-60}{3}=20$. La superfície màxima és
$A(30)=30\cdot 20=600\ \text{m}^2$.
\resposta{$x=30$ m, $y=20$ m; superfície màxima $600\ \text{m}^2$.}
\end{exresolt}

\begin{exresolt}[cost mínim per unitat]
El cost per unitat de servei (en euros) segons el nombre de centenars d'unitats $x$ és
$C(x)=x^2-24x+200$. Quantes unitats fan el cost mínim?
\solucio
$C'(x)=2x-24=0\Rightarrow x=12$. $C'$ passa de $-$ a $+$: \textbf{mínim}. Amb $x=12$
($1\,200$ unitats), el cost mínim per unitat és $C(12)=144-288+200=56\eur$.
\resposta{$x=12$ (1\,200 unitats); cost mínim $56\eur$ per unitat.}
\end{exresolt}

\subsection{Exercicis per fer}

\begin{exercicis}
\begin{exlist}
  \item Els ingressos d'una fira solidària (en euros) segons el preu d'entrada $x$ són
        $I(x)=-x^2+40x$. Troba el preu òptim i l'ingrés màxim.
  \item El cost d'un programa (en centenars d'euros) segons els mesos $x$ és
        $C(x)=x^2-50x+700$. Troba els mesos que el fan mínim i el cost mínim.
  \item Reparteix $18$ recursos en dues partides perquè el producte sigui màxim.
  \item Un hort rectangular ha de tenir un perímetre de $80$ m. Escriu l'àrea en funció
        d'un costat $x$, deriva i troba les mides de màxima superfície.
  \item \textbf{(Plantejament)} A partir d'una cartolina quadrada de $18$ cm de costat
        es fa una caixa sense tapa retallant quadrats de costat $x$. Escriu $V(x)$,
        deriva i troba el retall que dóna el volum màxim (descarta la solució sense
        sentit). \emph{(Pista: $V'(x)=12(x-3)(x-9)$.)}
  \item \textbf{(Interpretació)} En l'exemple de la satisfacció del casal, què vol dir,
        socialment, que passat el punt òptim la satisfacció baixi tot i tenir més
        participants? Com ho tindries en compte en planificar l'aforament?
\end{exlist}
\end{exercicis}

% ---------------------------------------------------------------------
\fullsolucions{Sessió 10}
\begin{solucions}
\begin{sollist}
  \item $I'(x)=-2x+40=0\Rightarrow x=20$ (màxim). Preu òptim $20\eur$; ingrés màxim
        $I(20)=-400+800=400\eur$.
  \item $C'(x)=2x-50=0\Rightarrow x=25$ (mínim). Cost mínim $C(25)=625-1250+700=75$
        centenars d'euros ($7\,500\eur$).
  \item $P(x)=x(18-x)$; $P'(x)=18-2x=0\Rightarrow x=9$ (màxim). Partides $9$ i $9$;
        producte $81$.
  \item Perímetre $2x+2y=80\Rightarrow y=40-x$; $A(x)=x(40-x)$;
        $A'(x)=40-2x=0\Rightarrow x=20$ (màxim). Costats $20\times 20$ m (quadrat); àrea
        $400\ \text{m}^2$.
  \item Base $18-2x$, altura $x$: $V(x)=x(18-2x)^2=4x^3-72x^2+324x$;
        $V'(x)=12x^2-144x+324=12(x-3)(x-9)$. Com que $x<9$, la solució vàlida és $x=3$
        (màxim). Volum màxim $V(3)=3\cdot 12^2=432\ \text{cm}^3$.
  \item Resposta oberta. Idea: passat l'òptim, afegir més participants \emph{empitjora}
        l'experiència (massificació, menys atenció per persona), de manera que «com més,
        millor» deixa de ser cert; en planificar l'aforament caldria fixar-lo a l'entorn
        del punt òptim, no al màxim físic possible.
\end{sollist}
\end{solucions}
